\subsection{The relative cylinder over an aspherical four-manifold}

Asphericity lets us turn a family of cylinder automorphisms into a
relative homotopy problem.  Surgery assembly then detects the normal
invariant of its mapping cylinder.  The resulting comparison separates
the assembly hypotheses needed for the component and fundamental-group
kernels from the additional hypothesis needed for component surjectivity.

Orson--Powell--Randal-Williams use realization of the centre by
automorphism loops for aspherical smooth four-manifolds in their
stabilized smooth pseudo-isotopy construction
\cite[Section~8.1, Remark~8.2 and Example~8.3, pp.~25--26;
Section~8.2, pp.~26--27]{OrsonPowellRandalWilliams2026}.
Here the same central evaluation condition enters the ordinary
PL-to-TOP comparison for the fully boundary-fixed cylinder, through
the relative structure classification and the assembly maps below.

Let $B$ be a closed connected oriented real PL four-manifold
which is aspherical, let $\pi=\pi_1B$, and put $N=B\times I$.
Use the ordinary TOP and joint PL family conventions of
Theorem~\ref{thm:four-dimensional-comparison-fibres}, and write
\[
 H_C=\operatorname{Aut}_C(N\ \mathrm{rel}\ \partial N),
 \qquad C\in\{\mathrm{PL},\mathrm{TOP}\}.
\]
Thus both ends are fixed pointwise.  Homotopy groups are
based at the identity; the TOP groups are replaced by
their ordinary singular complexes for comparison.
Suppose that $\operatorname{Wh}(\pi)=0$, and write
\begin{equation}\label{eq:aspherical-four-assembly}
 A_j:H_j(B;\mathbb L\langle1\rangle)
                   \longrightarrow L_j(\mathbb Z[\pi])
\end{equation}
for quadratic assembly with trivial orientation character
and the decoration of the geometric surgery sequence.
Whitehead-group vanishing identifies the simple and
homotopy decorations needed here.

\begin{theorem}\label{thm:aspherical-four-relative-comparison}
If $A_7$ is injective, the natural relative structure
boundary map identifies, as groups,
\begin{equation}\label{eq:aspherical-four-component-kernel}
 \ker\bigl(\pi_0H_{\mathrm{PL}}\longrightarrow
                  \pi_0H_{\mathrm{TOP}}\bigr)
                 \cong H^1(B;\mathbb Z/2).
\end{equation}
If $A_8$ is also injective, there is an exact sequence
\begin{equation}\label{eq:aspherical-four-fundamental-sequence}
 0\longrightarrow H^0(B;\mathbb Z/2)
 \longrightarrow\pi_1H_{\mathrm{PL}}
 \longrightarrow\pi_1H_{\mathrm{TOP}}
 \longrightarrow0,
\end{equation}
and the comparison is an isomorphism on $\pi_j$ for
every $j\geq2$.

If, in addition to $A_7$, the map $A_6$ is injective and
every element of $Z(\pi)$ is the evaluation class of a
based joint PL loop of automorphisms of $B$, then
\begin{equation}\label{eq:aspherical-four-component-sequence}
 0\longrightarrow H^1(B;\mathbb Z/2)
 \longrightarrow\pi_0H_{\mathrm{PL}}
 \longrightarrow\pi_0H_{\mathrm{TOP}}
 \longrightarrow0
\end{equation}
is exact as a sequence of groups.  Neither exact sequence
is asserted to split.

Let $\eta_B$ denote the ordinary-loop contribution to
\eqref{eq:four-dimensional-stabilization-components}:
\begin{equation}\label{eq:aspherical-four-ordinary-loop-map}
 \eta_B:\pi_1\operatorname{Aut}_{\mathrm{TOP}}(B)
 \longrightarrow\pi_0F_A
 \xrightarrow{\pi_0\beta_B}H^2(B;\mathbb Z/2).
\end{equation}
Under $A_6$-injectivity, its value on a loop is zero
whenever that loop's evaluation class is represented
by a based joint PL loop.  In particular the preceding
central evaluation hypothesis implies $\eta_B=0$.
That hypothesis is automatic if $Z(\pi)=0$.
\end{theorem}

\begin{proof}
All ordinary cohomology groups in this proof have
coefficients $\mathbb Z/2$.

\emph{A relative mapping-cylinder calculation.}
For $r\geq1$ put $V=B\times I^r$.  Suppose that a TOP
homeomorphism $F:V\to V$ fixes $\partial V$ and is
homotopic to the identity relative to that boundary.
Its stable PL reduction difference is
\[
 \kappa(F)\in H^3(V,\partial V)\cong H^{3-r}(B).
\]
We claim that
\begin{equation}\label{eq:aspherical-four-relative-zero-lemma}
 A_{5+r}\text{ injective}\quad\Longrightarrow\quad
 \kappa(F)=0.
\end{equation}

Choose a PL collar of the entire boundary of the
underlying PL manifold $V$, including the product corners
\cite[Theorem~2.25 and Corollary~2.26,
pp.~24--26]{RourkeSanderson1972}.
On its interval coordinate $0\leq t\leq\delta$, use
\[
 \sigma_s(t)=\min\{t+s,(t+\delta)/2\},
 \qquad 0\leq s\leq\delta/2,
\]
and use the identity off the collar.  This defines an
embedding $c_s:V\to V$ omitting the outer collar of
width $s$, with inverse collar coordinate
$\max\{t-s,2t-\delta\}$.  Conjugate $F$ by $c_s$ on
its image and extend by the identity on the omitted
collar.  The formulas agree at their common boundary;
the identical construction for $F^{-1}$ gives the inverse.
They define an isotopy relative to $\partial V$.
For a given relative homotopy of maps $k_u$, the formula
$c_sk_uc_s^{-1}$, likewise extended by the identity,
is a homotopy of maps with a common fixed collar.
Continuity follows by pasting on the two closed regions
in $(s,x,u)$; the displayed inverse collar coordinate
is continuous on the first region, including $s=0$.
This construction does not require $k_u$ to be invertible.

After this compression and an end-constant reparametrization,
take $k_0=1$ and $k_1=F$, both fixed on that collar.
The labeled mapping cylinder and its target are
\[
 M_F=(V\times I)\cup_FV,\qquad W=V\times I,
 \qquad U([x,u])=(k_u(x),u).
\]
At the upper face the label is $y=F(x)$, so $U$ is
the identity in that label.  It is also the identity
on the bottom and side faces.  The end-constant
homotopy and the fixed collar make these identifications
products near every face and corner intersection.
The bottom inclusions are homotopy equivalences and
$U$ restricts to their common identity.  Hence $U$ is
a homotopy equivalence of the full boundary pairs.

Its source stable TOP normal reduction, transported
with the Spivak normal-fibration comparison and the
specified boundary data, gives a relative normal
invariant satisfying
\begin{equation}\label{eq:aspherical-four-normal-invariant}
 \nu(U)\in[W,\partial W;G/\mathrm{TOP},*]
       \cong H_{5+r}(B;\mathbb L\langle1\rangle),
 \qquad \theta\nu(U)=0.
\end{equation}
The dimension here is $5+r\geq6$.  The compact relative
normal-invariant sequence and its identification with
assembly are precisely
\cite[Definition~18.2, pp.~197--198; relative-version
paragraph on p.~198; Proposition~18.3 and Theorem~18.5,
pp.~198--201]{Ranicki1992TopologicalManifolds}.
Thus injectivity of $A_{5+r}$ gives $\nu(U)=0$.

This zero class includes the fixed boundary comparison.
Forgetting the spherical-fibration comparison gives a
relative homotopy of the source TOP normal reduction
to the pullback of the target's standard reduction.
The latter has a PL lift.  Pulling this lift along the
homotopy gives a PL normal reduction of $M_F$ agreeing
with all its prescribed boundary reductions.  Its
relative Kirby--Siebenmann obstruction is consequently
zero.  The tangent and normal obstructions agree by
the Whitney-sum identity and stable triviality of their
sum.  The mapping-cylinder definition identifies this
class with the interval suspension of $\kappa(F)$.
Since
\[
 H^3(V,\partial V)\xrightarrow{\ \cong\ }
 H^4(V\times I,\partial(V\times I))
\]
is an isomorphism, the claim
\eqref{eq:aspherical-four-relative-zero-lemma} follows.

\emph{Relative homotopies for at least two interval factors.}
If $r\geq2$, the boundary $B\times\partial I^r$ is
connected and its fundamental group maps onto that
of $V$.  A homeomorphism fixed there induces the
identity on $\pi_1V$.  The full preimage of the boundary
in the universal cover is connected.  The lift fixing
one boundary point therefore fixes this entire preimage.
The universal cover of $V$ is contractible by asphericity.
A relative cellular construction, selecting one cell
in each deck orbit, gives a deck-equivariant homotopy
from that lift to the identity.  All extension
obstructions vanish in the contractible target.
It descends to a homotopy relative to the boundary.
The same argument fixes a closed collar whenever the
original map fixes that collar.

\emph{The family maps in the relative comparison.}
Let $\mathcal S_N$ be the sliced relative PL structure
space from the proof of
Theorem~\ref{thm:four-dimensional-comparison-fibres}.
Its construction uses the exterior-collared manifold
$B\times\mathbb R$, fixed outside the compact strip;
the classical orbit fibration therefore has an open
five-dimensional ambient manifold.  The relevant
identification is
\begin{equation}\label{eq:aspherical-four-structure-space}
 \mathcal S_N\simeq
 \operatorname{Map}_*(N/\partial N,K(\mathbb Z/2,3)),
 \qquad \pi_q\mathcal S_N=H^{2-q}(B).
\end{equation}
The orbit quotient contains the entire component of
the standard structure.  The orbit map sends a
homeomorphism to its transported standard structure.

Represent an element of $\pi_qH_{\mathrm{TOP}}$, for
$q=1,2$, by a based cube family $h_p$, $p\in I^q$.
The collar compression in the proof of the cited
comparison theorem fixes one common spatial collar
for the family.  Precomposing each parameter coordinate
by a PL function constant near its endpoints also
makes the family the identity on a parameter-boundary
collar.  Its trace
\begin{equation}\label{eq:aspherical-four-family-trace}
 \mathcal T_h:N\times I^q\longrightarrow N\times I^q,
 \qquad (x,p)\longmapsto(h_p(x),p)
\end{equation}
is an actual TOP homeomorphism fixed near the full
boundary.  Its inverse uses $h_p^{-1}$ in the same slice.

The image of the family under the orbit map is its
trace's stable PL difference in
\[
 H^3(N\times I^q,\partial(N\times I^q))
                 \cong H^{2-q}(B).
\]
Here the tangent data of an ordinary family can be
checked directly.  For a base point $(x,p)$ mapped
to $(h_p(x),p)$, send a nearby second point $(y,p')$ to
\[
 \bigl(h_{p+t(p'-p)}(y),p'\bigr),\qquad 0\leq t\leq1.
\]
Extend the family constantly outside the parameter
cube using its identity collar.  For every fixed base
point this is a homeomorphism germ, with inverse using
the inverse of the indicated slice.  It preserves
the zero section and the fixed boundary germs.
At $t=0$ it is the vertical tangent map with trivial
parameter directions; at $t=1$ it is the full tangent
map of the trace.  Their relative PL obstructions are
therefore equal.  This verifies the family-to-trace
identification with the specified ordinary TOP convention.

For the trace, $r=q+1\geq2$.  The preceding relative
homotopy argument and
\eqref{eq:aspherical-four-relative-zero-lemma} show that
\[
 \pi_1H_{\mathrm{TOP}}\longrightarrow H^1(B)
 \quad\text{vanishes under }A_7\text{-injectivity},
\]
and that
\[
 \pi_2H_{\mathrm{TOP}}\longrightarrow H^0(B)
 \quad\text{vanishes under }A_8\text{-injectivity}.
\]

\emph{Exactness and the group law.}
The comparison fibre is
$F_H\simeq\operatorname{Map}(B,K(\mathbb Z/2,1))$.
Consequently its exact sequence contains
\begin{equation}\label{eq:aspherical-four-relative-long-sequence}
\begin{aligned}
0&\longrightarrow\pi_2H_{\mathrm{PL}}
 \longrightarrow\pi_2H_{\mathrm{TOP}}
 \longrightarrow H^0(B)
 \longrightarrow\pi_1H_{\mathrm{PL}}\\
 &\longrightarrow\pi_1H_{\mathrm{TOP}}
 \longrightarrow H^1(B)
 \longrightarrow\pi_0H_{\mathrm{PL}}
 \longrightarrow\pi_0H_{\mathrm{TOP}}.
\end{aligned}
\end{equation}
The two zero maps give
\eqref{eq:aspherical-four-component-kernel},
\eqref{eq:aspherical-four-fundamental-sequence}, and
the degree-two isomorphism.  Since the fibre's higher
homotopy groups vanish, the higher isomorphisms follow
from the same sequence.

We verify the asserted group structure at the component
kernel.  The path-model fibre consists of a PL map $a$
and a TOP path from the identity to $a$.  Pointwise
composition of the paths makes its map to
$\pi_0H_{\mathrm{PL}}$ a group homomorphism.
The identification of its components with $H^1(B)$
uses the relative PL difference of those path traces,
whose boundary maps are PL.  The difference of a
composite is the sum of the differences with the usual
pullback action.

That action is trivial.  Every element of $H_C$ is
homotopic to the identity as an unrestricted map of
$N$, since it fixes the bottom.  Moreover the exact
sequence of the pair gives naturally
\[
 H^j(N,\partial N)\cong
 \operatorname{coker}\bigl(H^{j-1}(B)
       \xrightarrow{\Delta}H^{j-1}(B)^{\oplus2}\bigr).
\]
Both boundary restrictions are identities, so the
induced relative cohomology action is the identity.
For a path trace the same assertion follows after
the interval suspension: it preserves the parameter,
and acts identically on the relative cohomology of
every slice.  Equivalently, the relative product
filtration has only the single interval column.
The trace difference of a product therefore adds in
ordinary cohomology.  The identification with $H^1(B)$
and its injective boundary map in
\eqref{eq:aspherical-four-component-kernel} are group
homomorphisms, as claimed.

\emph{Components and central evaluation.}
For the single interval the boundary has two components.
If $F(x,t)=(a_t(x),v(x,t))$ fixes both ends, projecting
onto $B$ gives a loop $a_t$ in
$\operatorname{Map}(B,B)_{\mathrm{id}}$.  Each $a_t$
is homotopic to the identity and hence a homotopy
equivalence.  For an aspherical connected $B$,
the classical evaluation calculation gives
\cite[Lemma~4.3, p.~3336]{EnkelmannEtAl2018Atheory}
\begin{equation}\label{eq:aspherical-four-central-evaluation}
 \pi_1\operatorname{Map}(B,B)_{\mathrm{id}}\cong Z(\pi).
\end{equation}
Indeed a loop is a map $B\times S^1\to B$ restricting
to the identity on $B\times\{1\}$.  Its fundamental-group
homomorphism sends the extra circle generator to a
central element.  Conversely any such homomorphism
extends that identity and is realized uniquely up to
relative homotopy by the cellular $K(\pi,1)$ construction.
Loop concatenation corresponds to multiplication in
the centre.

Choose a joint PL automorphism loop $b_t$ representing
the same central element.  Its trace $T_b(x,t)=(b_t(x),t)$
is PL and fixes both ends.  The composition $T_b^{-1}F$
has first coordinate $b_{v(x,t)}^{-1}(a_t(x))$.
Interpolating $v(x,t)$ to $t$, both in this subscript
and in the output interval coordinate, is a homotopy
of maps relative to the two ends.  Its result is the
trace of the loop of maps $b_t^{-1}a_t$.  This loop
has zero central evaluation class by
\eqref{eq:aspherical-four-central-evaluation}, so is
nullhomotopic as a based loop of maps.  Thus
$T_b^{-1}F$ is homotopic to the identity relative to
both ends.  The relative zero lemma with $r=1$ and
$A_6$-injectivity gives $\kappa(T_b^{-1}F)=0$.
Since $T_b$ is PL and the relative cohomology action
is trivial, $\kappa(F)=0$.

In the five-dimensional relative structure space
\eqref{eq:aspherical-four-structure-space}, a zero
difference means membership in the standard component.
That entire component lies in the orbit quotient.
Its Kan fibration therefore lifts a structure path
to a TOP isotopy, with the prescribed fixed collar,
ending at a PL homeomorphism.  The preliminary collar
compression preserves the original relative component.
This proves component surjectivity and
\eqref{eq:aspherical-four-component-sequence}.

Finally, the trace identification in
Theorem~\ref{thm:four-dimensional-stabilization-fibre}
says that $\eta_B$ on an ordinary loop is the PL
difference of its trace.  Removing a joint PL loop
with the same evaluation class and applying the same
$r=1$ argument makes this difference zero.  Only
$A_6$-injectivity is needed for this assertion about
the specified loop.  In particular it vanishes on
the kernel of evaluation.  This proves the theorem.
\end{proof}

\begin{corollary}\label{cor:aspherical-four-explicit-groups}
All the hypotheses of
Theorem~\ref{thm:aspherical-four-relative-comparison}
hold for the standard $T^4$.  Its relative component
comparison is surjective with kernel $(\mathbb Z/2)^4$;
its fundamental-group comparison is surjective with
kernel $\mathbb Z/2$; its comparisons in degrees at
least two are isomorphisms; and $\eta_{T^4}=0$.

More generally, the Whitehead and assembly hypotheses
hold for every $B$ as above whose fundamental group
is word-hyperbolic or acts properly, cocompactly and
isometrically on a finite-dimensional CAT(0) space.
Thus the relative kernel, fundamental-group sequence
and higher isomorphisms hold throughout this class.
The component-surjectivity conclusion additionally
requires the stated realization of central elements
by joint PL loops, and holds in particular for
trivial centre.
\end{corollary}

\begin{proof}
For $T^4$, the free-abelian calculation
\cite[Example~24.16, p.~278]{Ranicki1992TopologicalManifolds}
gives $\operatorname{Wh}(\mathbb Z^4)=0$ and the
zero-connective quadratic assembly isomorphism in
degrees at least four.  The cofiber of
$\mathbb L\langle1\rangle\to\mathbb L\langle0\rangle$
is $H\mathbb Z$.  For $j=6,7,8$ the groups
$H_{j+1}(T^4;\mathbb Z)$ and $H_j(T^4;\mathbb Z)$
vanish, so the one-connective comparison in those
degrees is an isomorphism.  The translation loop
$b_t([x])=[x+tz]$ for $z\in\mathbb Z^4$ is jointly
PL and evaluates to $z$.  To see the joint assertion, work on
$[0,1]^4\times I$ and cut it by the finitely many affine hyperplanes
$x_i+t z_i=m$, with $m\in\mathbb Z$ meeting that compact domain.
On each resulting polytope the map modulo the lattice is affine.
A common triangulation, compatible with the quotient faces, makes
its trace jointly PL; the inverse uses $-z$ and the same construction.
The standard product
calculation gives $H^1(T^4;\mathbb Z/2)=(\mathbb Z/2)^4$
and $H^0(T^4;\mathbb Z/2)=\mathbb Z/2$.

For the general group class, first observe that $\pi$
is torsion-free.  The contractible universal cover
of the finite aspherical four-complex $B$ provides a
finite-length free $\mathbb Z\pi$ resolution of
$\mathbb Z$.  Restriction to a cyclic subgroup of
prime order would give a finite-length free resolution
over its group ring, contradicting its nonzero
cohomology in arbitrarily high degrees.  Every
nontrivial torsion element has a power of prime order.

Bartels--L\"uck's class contains both types of groups
required here
\cite[definition on p.~632]{BartelsLuck2012Borel}.
Their Theorem~B supplies coefficient K-theory
assembly isomorphisms in degrees at most zero,
surjectivity in degree one, and ultimate L-theory
assembly in every degree.  Proposition~0.3(i) and
its proof explicitly give, for torsion-free groups,
\begin{equation}\label{eq:aspherical-four-simple-assembly}
 \operatorname{Wh}(\pi)=0,\qquad
 H_j(B\pi;\mathbb L^s_{\mathbb Z})
       \xrightarrow{\ \cong\ }L_j^s(\mathbb Z\pi)
 \quad(j\in\mathbb Z)
\end{equation}
\cite[Theorem~B, pp.~633--634; Proposition~0.3 and
proof, pp.~634--635]{BartelsLuck2012Borel}.
The proof also gives vanishing of the reduced
projective class group and the negative K-groups.
Those are the specified inputs for the decoration
comparison and passage from virtually cyclic to
ordinary $B\pi$ assembly.  We use the trivial
orientation character and the algebraic conclusion
(i) of that proposition.

Finally, the cofiber $D$ of
$\mathbb L\langle1\rangle\to\mathbb L^s_{\mathbb Z}$
has no positive homotopy groups.  Filter the finite
four-complex $B$ by its cells.  A cell of dimension
$p\leq4$ contributes a suspension $\Sigma^pD$;
induction on the cells gives $H_j(B;D)=0$ for $j>4$.
The cofiber exact sequence, with outside terms
$H_{j+1}(B;D)$ and $H_j(B;D)$, therefore makes the
one-connective comparison an isomorphism for
$j=6,7,8$.  Equation~\eqref{eq:aspherical-four-simple-assembly}
proves all three required assembly isomorphisms.
\end{proof}

For a finite model of the torus, take the staircase
triangulation of the product of four three-edge
polygonal circles.  It has
\[
 (f_0,f_1,f_2,f_3,f_4)=(81,1215,4050,4860,1944).
\]
Indeed a $k$-simplex with minimal cubical carrier of
dimension $p$ corresponds to an ordered partition
of the $p$ varying coordinates into $k$ nonempty
blocks.  Thus, for $k\geq1$,
\[
 f_k=3^4\sum_{p=k}^4\binom4p\,k!\,S(p,k),
\]
where $S(p,k)$ is a Stirling number of the second kind.
The integral cellular chain complex is the tensor
product of four circle complexes, each chain homotopy
equivalent to one copy of $\mathbb Z$ in degrees zero
and one with zero differential.  Hence its integral
homology is free of ranks $(1,4,6,4,1)$.
This finite calculation describes the product model;
the assembly and automorphism comparisons follow from
the preceding geometric and algebraic arguments.

The conclusions evaluate the fully boundary-fixed
cylinder comparison and the ordinary-loop part of
$\beta_B$.  The remaining values of $\beta_B$, the
free-top comparison fibre, and the endpoint-image
condition require separate information.
